\documentclass[11pt,a4paper]{article}

\usepackage[utf8]{inputenc}
\usepackage[spanish,es-tabla]{babel}
\usepackage{amsmath,amssymb,amsfonts,amsthm}
\usepackage{geometry}
\usepackage{xcolor}
\usepackage{tcolorbox}
\usepackage{bm}

\geometry{top=2.5cm, bottom=2.5cm, left=2.5cm, right=2.5cm}

\tcbuselibrary{skins,breakable}

\newtcolorbox{solbox}[1]{
  colback=blue!5!white,
  colframe=blue!65!black,
  fonttitle=\bfseries,
  title=#1,
  arc=2mm,
  breakable
}

\newtcolorbox{resumenbox}{
  colback=green!5!white,
  colframe=green!50!black,
  fonttitle=\bfseries,
  title=Resumen de Resultados,
  arc=2mm
}

\title{\textbf{Geometría Lineal y Afín}\\ \large Resolución Formal: Planos Afines, Pertenencia y Sistemas de Coordenadas}
\author{Curso de Postgrado en Matemáticas}
\date{\today}

\begin{document}

\maketitle

\section*{Planteamiento del Problema}

En el espacio afín tridimensional real $\mathbb{A}^3(\mathbb{R})$, se consideran los siguientes puntos:
\[
P_0 = (1, 2, 0), \quad P_1 = (3, 2, 1), \quad P_2 = (0, 4, 1)
\]

\begin{enumerate}
    \item Determine las ecuaciones paramétricas e implícitas del plano afín $\Pi$ que contiene a $\{P_0, P_1, P_2\}$.
    \item Verifique si el punto $Q = (4, 4, 3)$ pertenece al plano afín $\Pi$.
    \item En caso afirmativo, calcule las coordenadas afines de $Q$ respecto al marco afín $\{P_0; v_1, v_2\}$ y determine sus coordenadas baricéntricas con respecto al símplice de referencia $\{P_0, P_1, P_2\}$.
\end{enumerate}

\hrule
\vspace{0.4cm}

\section{Comprobación Preliminar: Independencia Afín}

Recordemos que tres puntos $\{P_0, P_1, P_2\}$ en un espacio afín son afínmente independientes (no colineales) si y solo si los vectores directores con origen común:
\[
v_1 = \vec{P_0 P_1} = P_1 - P_0, \qquad v_2 = \vec{P_0 P_2} = P_2 - P_0
\]
son linealmente independientes en el espacio vectorial asociado $\mathbb{R}^3$.

Calculamos las componentes:
\begin{align*}
    v_1 &= (3 - 1,\, 2 - 2,\, 1 - 0) = (2, 0, 1), \\
    v_2 &= (0 - 1,\, 4 - 2,\, 1 - 0) = (-1, 2, 1).
\end{align*}

Analizando el rango de la matriz cuyas filas son $v_1$ y $v_2$:
\[
\operatorname{rg} \begin{pmatrix} 2 & 0 & 1 \\ -1 & 2 & 1 \end{pmatrix} = 2
\]
puesto que el menor principal de orden 2:
\[
\begin{vmatrix} 2 & 0 \\ -1 & 2 \end{vmatrix} = 4 \neq 0.
\]
Por lo tanto, $\{v_1, v_2\}$ es libre, lo cual demuestra que los puntos son afínmente independientes y determinan unívocamente un plano afín bidimensional $\Pi = P_0 + \langle v_1, v_2 \rangle$.

\section{Ecuaciones del Plano Afín $\Pi$}

\subsection{Ecuaciones Paramétricas}
Un punto arbitrario $X = (x, y, z) \in \mathbb{A}^3(\mathbb{R})$ pertenece a $\Pi$ si y solo si existen escalares $\lambda, \mu \in \mathbb{R}$ tales que:
\[
X = P_0 + \lambda v_1 + \mu v_2
\]
Sustituyendo los valores conocidos:
\[
\begin{pmatrix} x \\ y \\ z \end{pmatrix} = \begin{pmatrix} 1 \\ 2 \\ 0 \end{pmatrix} + \lambda \begin{pmatrix} 2 \\ 0 \\ 1 \end{pmatrix} + \mu \begin{pmatrix} -1 \\ 2 \\ 1 \end{pmatrix}
\]

Desglosando en componentes escalares obtenemos el sistema de ecuaciones paramétricas:
\begin{solbox}{Ecuaciones Paramétricas de $\Pi$}
\[
\Pi:\ \begin{cases}
x = 1 + 2\lambda - \mu \\
y = 2 + 2\mu \\
z = \lambda + \mu
\end{cases} \quad (\lambda, \mu \in \mathbb{R})
\]
\end{solbox}

\subsection{Ecuación Implícita (Cartesiana)}
La condición $X \in \Pi$ equivale a que el vector $\vec{P_0 X} = (x-1, y-2, z)$ pertenezca al subespacio vectorial director $\vec{\Pi} = \langle v_1, v_2 \rangle$. En dimensión 3, esto equivale a la anulación del determinante formado por estos tres vectores:
\[
\det \begin{pmatrix}
x - 1 & y - 2 & z \\
2 & 0 & 1 \\
-1 & 2 & 1
\end{pmatrix} = 0
\]

Desarrollando por la primera fila:
\[
(x - 1) \begin{vmatrix} 0 & 1 \\ 2 & 1 \end{vmatrix} - (y - 2) \begin{vmatrix} 2 & 1 \\ -1 & 1 \end{vmatrix} + z \begin{vmatrix} 2 & 0 \\ -1 & 2 \end{vmatrix} = 0
\]
\[
-2(x - 1) - 3(y - 2) + 4z = 0
\]
\[
-2x + 2 - 3y + 6 + 4z = 0 \implies -2x - 3y + 4z + 8 = 0
\]
Multiplicando por $-1$:
\begin{solbox}{Ecuación Implícita de $\Pi$}
\[
\Pi:\ 2x + 3y - 4z - 8 = 0
\]
\end{solbox}

\section{Verificación de Pertenencia del Punto $Q = (4, 4, 3)$}

Para determinar si $Q \in \Pi$, evaluamos directamente sus coordenadas en la ecuación implícita:
\[
2(4) + 3(4) - 4(3) - 8 = 8 + 12 - 12 - 8 = 20 - 20 = 0
\]
La ecuación se satisface de forma idéntica ($0 = 0$).

\begin{solbox}{Conclusión de Pertenencia}
El punto $Q = (4, 4, 3)$ satisface la ecuación implícita del plano, por lo tanto:
\[
\bm{Q \in \Pi} \quad \text{(Afirmativo)}
\]
\end{solbox}

\section{Sistemas de Coordenadas de $Q$}

\subsection{Coordenadas Afines respecto al marco afín $\mathcal{R} = \{P_0; v_1, v_2\}$}
Por definición de coordenadas afines en $\Pi$, las coordenadas $(\alpha_1, \alpha_2)$ de $Q$ son los únicos coeficientes de la descomposición del vector $\vec{P_0 Q}$ en la base del subespacio director $\{v_1, v_2\}$:
\[
\vec{P_0 Q} = \alpha_1 v_1 + \alpha_2 v_2
\]

Calculamos el vector $\vec{P_0 Q}$:
\[
\vec{P_0 Q} = Q - P_0 = (4 - 1,\, 4 - 2,\, 3 - 0) = (3, 2, 3)
\]

Planteamos el sistema matricial:
\[
\alpha_1 \begin{pmatrix} 2 \\ 0 \\ 1 \end{pmatrix} + \alpha_2 \begin{pmatrix} -1 \\ 2 \\ 1 \end{pmatrix} = \begin{pmatrix} 3 \\ 2 \\ 3 \end{pmatrix} \iff \begin{cases}
2\alpha_1 - \alpha_2 = 3 & (1) \\
2\alpha_2 = 2 & (2) \\
\alpha_1 + \alpha_2 = 3 & (3)
\end{cases}
\]

De la ecuación (2):
\[
\alpha_2 = 1
\]
Sustituyendo $\alpha_2 = 1$ en la ecuación (3):
\[
\alpha_1 + 1 = 3 \implies \alpha_1 = 2
\]
Comprobamos en la ecuación (1):
\[
2(2) - (1) = 4 - 1 = 3 \quad \checkmark
\]

\begin{solbox}{Coordenadas Afines}
Las coordenadas afines de $Q$ respecto al marco afín $\mathcal{R} = \{P_0; v_1, v_2\}$ son:
\[
[Q]_{\mathcal{R}} = (\alpha_1, \alpha_2) = \bm{(2, 1)}
\]
\end{solbox}

\subsection{Coordenadas Baricéntricas respecto al símplice $\{P_0, P_1, P_2\}$}
Dado el 2-símplice afín de referencia formado por los vértices $\{P_0, P_1, P_2\}$, las coordenadas baricéntricas normalizadas $(\lambda_0, \lambda_1, \lambda_2)$ se definen mediante:
\[
Q = \lambda_0 P_0 + \lambda_1 P_1 + \lambda_2 P_2, \quad \text{con} \quad \lambda_0 + \lambda_1 + \lambda_2 = 1
\]

Utilizando la relación con las coordenadas afines:
\[
Q = P_0 + \vec{P_0 Q} = P_0 + \alpha_1(P_1 - P_0) + \alpha_2(P_2 - P_0)
\]
\[
Q = (1 - \alpha_1 - \alpha_2) P_0 + \alpha_1 P_1 + \alpha_2 P_2
\]

Igualando coeficientes:
\begin{align*}
    \lambda_1 &= \alpha_1 = 2 \\
    \lambda_2 &= \alpha_2 = 1 \\
    \lambda_0 &= 1 - \alpha_1 - \alpha_2 = 1 - 2 - 1 = -2
\end{align*}

Comprobación de la suma de pesos:
\[
\lambda_0 + \lambda_1 + \lambda_2 = -2 + 2 + 1 = 1 \quad \checkmark
\]

Comprobación directa de la combinación afín:
\[
-2(1, 2, 0) + 2(3, 2, 1) + 1(0, 4, 1) = (-2+6+0,\, -4+4+4,\, 0+2+1) = (4, 4, 3) = Q \quad \checkmark
\]

\begin{solbox}{Coordenadas Baricéntricas}
Las coordenadas baricéntricas de $Q$ respecto al símplice $\{P_0, P_1, P_2\}$ son:
\[
(\lambda_0, \lambda_1, \lambda_2) = \bm{(-2, 2, 1)}
\]
\end{solbox}

\paragraph{Interpretación geométrica:} Al ser $\lambda_0 = -2 < 0$, el punto $Q$ se encuentra en el exterior de la envolvente convexa del símplice (fuera del triángulo afín $\triangle P_0 P_1 P_2$), situado del lado opuesto de la recta afín soporte generada por $P_1$ y $P_2$.

\vspace{0.4cm}

\begin{resumenbox}
\begin{itemize}
    \item \textbf{Ecuaciones paramétricas:} $x = 1 + 2\lambda - \mu$, \quad $y = 2 + 2\mu$, \quad $z = \lambda + \mu$.
    \item \textbf{Ecuación implícita:} $2x + 3y - 4z - 8 = 0$.
    \item \textbf{Pertenencia:} $Q \in \Pi$ (Verdadero).
    \item \textbf{Coordenadas afines en $\{P_0; v_1, v_2\}$:} $[Q] = (2, 1)$.
    \item \textbf{Coordenadas baricéntricas en $\{P_0, P_1, P_2\}$:} $(\lambda_0, \lambda_1, \lambda_2) = (-2, 2, 1)$.
\end{itemize}
\end{resumenbox}

\end{document}

